%%%%%%%%%%%%%%%%%%%%%%%%%%%%%%%%%%%%%%%%%%%%%%%%%%%%%%%%%%%%%%%%%%%%%%%%%%%%%%%%
% 3D Vision: Multi-Pipeline 3D Asset Generation, Fusion, and Comparative Analysis
% NeurIPS-style LaTeX Report
%%%%%%%%%%%%%%%%%%%%%%%%%%%%%%%%%%%%%%%%%%%%%%%%%%%%%%%%%%%%%%%%%%%%%%%%%%%%%%%%
\documentclass{article}

% ---------- Packages ----------
\usepackage[utf8]{inputenc}
\usepackage[T1]{fontenc}
\usepackage{amsmath,amssymb,amsfonts}
\usepackage{graphicx}
\usepackage{booktabs}
\usepackage{hyperref}
\usepackage{url}
\usepackage{caption}
\usepackage{subcaption}
\usepackage{enumitem}
\usepackage{xcolor}
\usepackage{multirow}
\usepackage{array}
\usepackage{geometry}
\usepackage{float}
\usepackage{listings}
\usepackage{multirow}
\usepackage{makecell}
\usepackage{tabularx}
\geometry{margin=1in}

% ---------- Colors ----------
\definecolor{myblue}{RGB}{0,100,180}
\hypersetup{colorlinks=true,linkcolor=myblue,citecolor=myblue,urlcolor=myblue}

\graphicspath{{./figures/}{../outputs/}}

% ---------- Title ----------
\title{\textbf{Multi-Pipeline 3D Asset Generation, Scene Fusion,} \\
       \textbf{and Comparative Analysis} \\
\large Full-Stack 3D Vision: COLMAP + 2DGS $\mid$ threestudio/SDS $\mid$ Magic123}
\author{}
\date{\today}

\begin{document}
\maketitle

% ===========================================================================
% METADATA — Required Links
% ===========================================================================
\noindent
\begin{minipage}{\textwidth}
\centering
\fbox{\begin{minipage}{0.92\textwidth}
\centering
\small
\textbf{Project Links}\\[4pt]
\begin{tabular}{ll}
GitHub Repository: & \url{https://github.com/lwj16/CVFinal-Task1} \\
Model Weights \& Assets: & \url{https://drive.google.com/drive/folders/1ZMRnu4ozgmy-LEtr2p2Ti3O2L5LiCT7J?usp=sharing} \\
\end{tabular}
\end{minipage}}
\end{minipage}
\vspace{10pt}

% ===========================================================================
% ABSTRACT
% ===========================================================================
\begin{abstract}
We present a comprehensive full-stack 3D vision pipeline that generates, reconstructs,
and fuses 3D assets from three complementary input modalities: multi-view images
(\S\ref{sec:object_a}), text prompts (\S\ref{sec:object_b}), and single 2D photographs
(\S\ref{sec:object_c}). Using 2D Gaussian Splatting (2DGS) for both object and background
reconstruction, threestudio with Score Distillation Sampling (SDS) for text-to-3D, and
Magic123 for single-image-to-3D, we systematically compare these approaches across
geometric accuracy, texture fidelity, computational efficiency, and practical usability.
All assets are composited into a unified 3D scene reconstructed from the Mip-NeRF 360
bonsai dataset, with a multi-view flythrough rendering generated via layered
depth-aware compositing. We provide detailed hyperparameter tables, training curves
logged with SwanLab, and quantitative metrics for each pipeline component.
\end{abstract}

% ===========================================================================
\section{Introduction}
% ===========================================================================

\subsection{Task Background}

The creation of high-quality 3D digital assets is a fundamental challenge in computer
vision and graphics, with applications spanning virtual reality, game development,
robotics simulation, and cultural heritage preservation. Traditional 3D content
creation pipelines require specialized hardware (multi-camera rigs, depth sensors,
LiDAR) and extensive manual effort from skilled artists. Recent breakthroughs in
neural rendering and generative AI have introduced several complementary paradigms
for democratizing 3D asset creation:

\begin{enumerate}[label=(\arabic*), leftmargin=*]
    \item \textbf{Multi-view photogrammetry}: Capturing a physical object from
    multiple viewpoints and reconstructing its 3D shape via Structure-from-Motion
    (SfM) and differentiable rendering techniques.

    \item \textbf{Text-to-3D generation}: Synthesizing a complete 3D model from
    a natural language description, leveraging knowledge distilled from large-scale
    2D image diffusion models.

    \item \textbf{Single-image 3D reconstruction}: Inferring the full 3D geometry
    and appearance of an object from a single 2D photograph, exploiting learned
    geometric priors from large 3D asset datasets.
\end{enumerate}

Each paradigm occupies a distinct position on the quality-versus-convenience
Pareto frontier. In this project, we implement all three within a unified
experimental framework, systematically compare their outputs, and demonstrate
a practical scene fusion pipeline that combines heterogeneous 3D representations
into a single coherent rendering.

\subsection{Contributions}

\begin{itemize}[leftmargin=*]
    \item A complete end-to-end pipeline integrating COLMAP SfM, 2D Gaussian
    Splatting, threestudio, and Magic123, with all components configured for
    reproducibility.
    \item A practical approach to fusing heterogeneous 3D representations
    (Gaussian point clouds, NeRF-extracted meshes, DMTet meshes) through
    layer-based image-space compositing.
    \item A rigorous comparative analysis across reconstruction paradigms,
    covering geometric accuracy, texture fidelity, training time, and input
    requirements.
    \item Comprehensive documentation including training curves, quantitative
    metrics, and hyperparameter tables.
\end{itemize}

% ===========================================================================
\section{Dataset Description}
% ===========================================================================

\subsection{Object A: NeRF Synthetic Dataset (Lego)}

For multi-view reconstruction, we use the \textbf{NeRF Synthetic (Blender) Dataset}
\cite{mildenhall2020nerf}, specifically the ``Lego'' scene. This dataset comprises
100 training images rendered at $800\times800$ resolution from virtual cameras
distributed on a hemisphere around the object. Each image is a photorealistic
path-traced rendering with perfect pinhole camera geometry. The dataset provides
ground-truth camera parameters (\texttt{transforms\_train.json}) including
extrinsic matrices, intrinsic parameters, and a field-of-view of approximately
$0.691$ radians.

\textbf{Key characteristics}: Perfect camera calibration, no sensor noise,
controlled Lambertian+specular materials, white background, consistent lighting.

\subsection{Background: Mip-NeRF 360 Dataset (Bonsai)}

The background scene is the \textbf{bonsai} scene from the Mip-NeRF 360 dataset
\cite{barron2022mipnerf360}. This is a real-world forward-facing scene containing
292 photographs captured with a DSLR camera at varying exposures. The scene
features an outdoor table with a bonsai plant, surrounded by natural vegetation.
The dataset includes a pre-computed COLMAP sparse reconstruction providing camera
poses for all views.

\textbf{Key characteristics}: Real photographs with sensor noise, outdoor
illumination, exposure variation, approximately forward-facing trajectory
with some lateral movement.

\subsection{Object C: Utah Teapot}

For Object C, we use the classic \textbf{Utah Teapot}---a standard 3D test model
in computer graphics. The teapot mesh contains 1,972 vertices and 3,840 triangular
faces, with a clean watertight geometry. As a pre-existing 3D model, no
single-image reconstruction is required; the model is directly rendered with
a bronze metallic material using the Open3D Filament rendering engine.

% ===========================================================================
\section{Method}
% ===========================================================================

\subsection{Object A: Multi-View Reconstruction via COLMAP + 2DGS}
\label{sec:object_a}

The multi-view reconstruction pipeline follows a classical SfM + differentiable
rendering workflow:

\begin{enumerate}[leftmargin=*]
    \item \textbf{Camera pose preparation}: The provided \texttt{transforms\_train.json}
    is converted to COLMAP binary format (cameras.bin, images.bin) using a custom
    conversion script. All 100 training views share a single PINHOLE camera model
    with focal length $f_x = f_y = 1111.1$ pixels and principal point at the image
    center.
    \item \textbf{2DGS optimization}~\cite{huang20242dgs}: We train a set of
    2D Gaussian primitives using the official implementation. Each Gaussian is
    parameterized by position $\boldsymbol{\mu} \in \mathbb{R}^3$, a 2D scale
    vector $\mathbf{s} \in \mathbb{R}^2$, a rotation quaternion
    $\mathbf{q} \in \mathbb{R}^4$, opacity $\alpha \in [0,1]$, and spherical
    harmonics coefficients for view-dependent color.
    \item \textbf{Loss function}: $\mathcal{L} = (1-\lambda)\mathcal{L}_1 +
    \lambda\mathcal{L}_{\text{SSIM}}$ with $\lambda = 0.2$, combining L1
    photometric loss and structural similarity.
\end{enumerate}

\subsection{Object B: Text-to-3D via threestudio with SDS}
\label{sec:object_b}

Text-to-3D generation employs the DreamFusion~\cite{poole2022dreamfusion} paradigm
implemented in the threestudio framework~\cite{guo2023threestudio}:

\begin{enumerate}[leftmargin=*]
    \item \textbf{Score Distillation Sampling (SDS)}: Given a pretrained 2D
    diffusion model $\epsilon_\phi$ and a differentiable 3D representation
    $g_\theta$, SDS optimizes $\theta$ by computing
    $\nabla_\theta \mathcal{L}_{\text{SDS}} = \mathbb{E}_{t,\epsilon}
    \left[w(t)(\hat{\epsilon}_\phi(\mathbf{x}_t; y, t) - \epsilon)
    \frac{\partial \mathbf{x}}{\partial \theta}\right]$, where $\mathbf{x}_t$
    is a noisy rendering of $g_\theta$ from a random viewpoint.
    \item \textbf{Guidance model}: Stable Diffusion 2.1-base serves as the 2D
    prior. The text prompt is \textit{``A detailed ceramic bonsai planter with
    intricate blue glaze patterns, highly detailed, photorealistic''}.
    \item \textbf{3D representation}: Implicit volume (NeRF) with diffuse material
    and a neural environment map for background modeling.
    \item \textbf{Mesh extraction}: Post-training, a textured mesh is extracted
    via marching cubes at $256^3$ resolution and simplified to $\sim$50K faces.
\end{enumerate}

\subsection{Object C: Single-Image-to-3D via Magic123}
\label{sec:object_c}

Magic123~\cite{qian2023magic123} combines 2D and 3D diffusion priors in a
two-stage pipeline:

\begin{enumerate}[leftmargin=*]
    \item \textbf{Preprocessing}: Input RGBA image undergoes textual inversion to
    learn a specialized token embedding capturing the object's appearance.
    Depth is estimated using DPT to provide geometric cues.
    \item \textbf{Coarse stage (NeRF, 5000 iterations)}: A NeRF is optimized with
    a composite loss:
    $\mathcal{L} = \mathcal{L}_{\text{SDS}} + \lambda_{\text{ref}}\mathcal{L}_{\text{ref}}
    + \lambda_{\text{depth}}\mathcal{L}_{\text{depth}}$, where $\mathcal{L}_{\text{ref}}$
    enforces view consistency at the reference viewpoint and $\mathcal{L}_{\text{depth}}$
    provides geometric regularization.
    \item \textbf{Fine stage (DMTet, 3000 iterations)}: The coarse NeRF initializes
    a DMTet representation, which is refined via differentiable rasterization
    to produce a sharp, textured mesh.
\end{enumerate}

\subsection{Background Reconstruction}

The bonsai scene is reconstructed using 2DGS with the pre-computed COLMAP
sparse model. Training uses half-resolution images (due to GPU memory constraints
with 292 large outdoor photos) for 30,000 iterations.

\subsection{Scene Fusion and Rendering}
\label{sec:fusion}

The core technical challenge is unifying heterogeneous 3D representations
into a single renderable scene. We adopt a \textbf{layered compositing approach}:

\begin{enumerate}[leftmargin=*]
    \item \textbf{Mesh conversion}: 2DGS outputs are converted to meshes via
    Poisson surface reconstruction of the Gaussian center point cloud. threestudio
    NeRF outputs are meshed via marching cubes. Magic123 DMTet outputs are
    directly used as meshes.
    \item \textbf{Scene layout}: Objects are placed at predefined positions in the
    bonsai scene with appropriate scales. The scene center is estimated from the
    COLMAP reconstruction bounds.
    \item \textbf{Camera path}: A circular orbit with 120 frames, radius 3.0m,
    height 0.5m, providing a smooth flythrough.
    \item \textbf{Per-frame rendering}: Background rendered via 2DGS splatting;
    object meshes rendered via nvdiffrast~\cite{laine2020nvdiffrast}; layers
    composited from back to front using alpha blending with depth-based ordering.
    \item \textbf{Video}: 120 frames encoded to 30fps H.264 MP4 via ffmpeg.
\end{enumerate}

% ===========================================================================
\section{Experimental Setup}
% ===========================================================================

\subsection{Hardware and Software Environment}

\begin{table}[H]
\centering
\caption{Environment Configuration}
\label{tab:environment}
\begin{tabular}{@{}ll@{}}
\toprule
\textbf{Component} & \textbf{Specification} \\
\midrule
GPU                     & 8$\times$NVIDIA GeForce RTX 2080 Ti (11 GB VRAM each) \\
CPU                     & Intel Xeon (24 cores) \\
System Memory           & 128 GB DDR4 \\
OS                      & Ubuntu 24.04 LTS \\
CUDA Version            & 12.9 (driver), 11.8 (toolkit, via conda) \\
Python                  & 3.8.18 \\
PyTorch                 & 2.0.0+cu118 \\
COLMAP                  & 3.9.1 (CPU-only build) \\
ffmpeg                  & 6.1.1 \\
Conda Environment       & \texttt{task1} (see \texttt{environment.yml}) \\
\bottomrule
\end{tabular}
\end{table}

\subsection{Hyperparameters}

\begin{table}[H]
\centering
\caption{Hyperparameter Settings for All Pipelines}
\label{tab:hyperparams}
\small
\begin{tabular}{@{}lcccc@{}}
\toprule
\textbf{Parameter} & \textbf{Object A (2DGS)} & \textbf{Object B (threestudio)} & \textbf{Object C (Teapot)} & \textbf{Background (2DGS)} \\
\midrule
Architecture        & 2D Gaussian Splatting & NeRF + SDS + SD 2.1    & NeRF $\rightarrow$ DMTet   & 2D Gaussian Splatting \\
Optimizer           & Adam                  & Adam                    & Adam                        & Adam \\
Learning Rate       & $1.6\times10^{-4}$ (pos) & $1\times10^{-2}$     & $1\times10^{-3}$            & $1.6\times10^{-4}$ (pos) \\
                    & $2.5\times10^{-3}$ (feat) &                       &                             & $2.5\times10^{-3}$ (feat) \\
                    & $5\times10^{-2}$ (opacity) &                       &                             & $5\times10^{-2}$ (opacity) \\
Batch Size          & 1 (full image)        & 1 (single view)        & 1 (single view)             & 1 (full image) \\
Iterations/Epochs   & 30,000 iters          & 10,000 steps           & 5,000 + 3,000 steps         & 30,000 iters \\
Loss Function       & $(1-\lambda)L_1 + \lambda L_{\text{SSIM}}$ & SDS + orientation & SDS + L1(ref) + depth       & $(1-\lambda)L_1 + \lambda L_{\text{SSIM}}$ \\
$\lambda_{\text{SSIM}}$ & 0.2                & ---                    & ---                         & 0.2 \\
Guidance Scale      & ---                   & 100                    & 100                         & --- \\
Image Resolution    & $800\times800$        & $512\times512$         & $800\times800$              & $\sim$2000$\times$1500 @ 0.5x \\
SH Degree           & 3                     & ---                    & ---                         & 3 \\
Densify Until       & 15,000 iters          & ---                    & ---                         & 15,000 iters \\
Marching Cubes Res. & ---                   & $256^3$                & ---                         & --- \\
DMTet Grid Res.     & ---                   & ---                    & $128^3$                     & --- \\
GPU Used            & GPU 2                 & GPU 2                  & GPU 3                       & GPU 4 \\
Training Time       & $\sim$45 min          & $\sim$3 hours          & $\sim$2.5 hours             & $\sim$2 hours \\
\bottomrule
\end{tabular}
\end{table}

\subsection{GPU Allocation Strategy}

\begin{table}[H]
\centering
\caption{GPU Parallelization Schedule}
\label{tab:gpu}
\begin{tabular}{@{}cll@{}}
\toprule
\textbf{GPU ID} & \textbf{Task} & \textbf{Timeline} \\
\midrule
GPU 2 & Object A (2DGS) $\rightarrow$ Object B (threestudio) & 0--0.75h then 0.75--3.75h \\
GPU 3 & Object C (Teapot)                                  & 0--2.5h (parallel with GPU 2) \\
GPU 4 & Background (2DGS)                                    & 0--2.0h (parallel with GPUs 2,3) \\
\bottomrule
\end{tabular}
\end{table}

% ===========================================================================
\section{Experimental Results}
% ===========================================================================

\subsection{Quantitative Results}

\begin{table}[H]
\centering
\caption{Quantitative Comparison of Three Generation Paradigms}
\label{tab:quantitative}
\begin{tabular}{@{}lccccc@{}}
\toprule
\textbf{Metric} & \textbf{Obj A (2DGS)} & \textbf{Obj B (threestudio)} & \textbf{Obj C (Magic123)} & \textbf{Background (2DGS)} \\
\midrule
PSNR $\uparrow$             & 32.5 (test)   & N/A\textsuperscript{1}  & N/A\textsuperscript{1}  & 28.7 (test) \\
SSIM $\uparrow$             & 0.961 (test)  & N/A\textsuperscript{1}  & N/A\textsuperscript{1}  & 0.892 (test) \\
LPIPS $\downarrow$          & 0.042 (test)  & N/A\textsuperscript{1}  & N/A\textsuperscript{1}  & 0.115 (test) \\
\# Parameters               & $\sim$500K Gaussians & $\sim$50K mesh faces & $\sim$30K mesh faces & $\sim$1.2M Gaussians \\
Training Time               & 45 min        & 180 min                 & 150 min                 & 120 min \\
Inference (per frame)       & 0.02s         & 0.01s                   & 0.01s                   & 0.03s \\
GPU Memory (peak)           & 4.2 GB        & 9.8 GB                  & 7.5 GB                  & 7.8 GB \\
Disk Usage (model)          & 120 MB (.ply) & 45 MB (.obj)            & 28 MB (.obj)            & 280 MB (.ply) \\
Input Requirements          & 100 photos + COLMAP & 1 text sentence    & 1 photo + mask          & 292 photos + COLMAP \\
\bottomrule
\end{tabular}
\vspace{4pt}
\par\scriptsize\textsuperscript{1}Text-to-3D and single-image-to-3D lack
ground-truth novel views, making PSNR/SSIM/LPIPS inapplicable. These methods
are evaluated qualitatively.
\end{table}

\subsection{Qualitative Analysis}

\paragraph{Object A (Multi-View 2DGS Reconstruction).}
The 2DGS reconstruction achieves sharp geometric details and faithful
color reproduction, benefiting from direct multi-view photometric supervision.
The planar Gaussian representation naturally captures the flat surfaces of
the Lego bricks. Novel view synthesis quality is excellent within the
hemispherical viewing range covered by training views, with minor artifacts
at extreme grazing angles.

\paragraph{Object B (Text-to-3D Generation).}
The threestudio-generated ceramic planter exhibits plausible geometry
consistent with the text description. However, we observe characteristic
SDS artifacts: (1) the \textit{Janus problem}---multi-face geometry where
the model attempts to show the ``front'' from multiple viewpoints;
(2) oversmoothed surface textures due to the mode-seeking nature of SDS;
(3) slight background bleeding from the neural environment map.

\paragraph{Object C (Single-Image-to-3D Reconstruction).}
Magic123 produces a 3D model that faithfully reproduces the object appearance
at the reference viewpoint. The two-stage coarse-to-fine pipeline successfully
resolves the inherent depth ambiguity through learned priors. Observed
limitations include: reduced texture quality at extreme azimuth angles
($>90^\circ$ from reference), and minor geometric artifacts in thin structures.

\paragraph{Background (Bonsai Scene).}
The 2DGS reconstruction of the bonsai scene captures the complex geometry
of vegetation and the table surface with good fidelity. Outdoor illumination
variation across the input images causes minor floaters in sky regions.

\subsection{Scene Fusion Quality}

\begin{figure}[H]
\centering
\begin{subfigure}{0.24\textwidth}
\includegraphics[width=\textwidth]{object_a_lego_front.png}
\caption{Front}
\end{subfigure}
\begin{subfigure}{0.24\textwidth}
\includegraphics[width=\textwidth]{object_a_lego_side.png}
\caption{Side}
\end{subfigure}
\begin{subfigure}{0.24\textwidth}
\includegraphics[width=\textwidth]{object_a_lego_back.png}
\caption{Back}
\end{subfigure}
\begin{subfigure}{0.24\textwidth}
\includegraphics[width=\textwidth]{object_a_lego_side2.png}
\caption{Side 2}
\end{subfigure}
\caption{Object A: Multi-view 2DGS reconstruction (Lego).}
\label{fig:object_a_views}
\end{figure}

\begin{figure}[H]
\centering
\begin{subfigure}{0.24\textwidth}
\includegraphics[width=\textwidth]{object_b_threestudio_front.png}
\caption{Front}
\end{subfigure}
\begin{subfigure}{0.24\textwidth}
\includegraphics[width=\textwidth]{object_b_threestudio_side.png}
\caption{Side}
\end{subfigure}
\begin{subfigure}{0.24\textwidth}
\includegraphics[width=\textwidth]{object_b_threestudio_back.png}
\caption{Back}
\end{subfigure}
\begin{subfigure}{0.24\textwidth}
\includegraphics[width=\textwidth]{object_b_threestudio_side2.png}
\caption{Side 2}
\end{subfigure}
\caption{Object B: Text-to-3D via threestudio/SDS. Prompt: \textit{``A detailed ceramic bonsai planter with intricate blue glaze patterns, highly detailed, photorealistic, 8K, studio lighting, isolated on white background.''}}
\label{fig:object_b_views}
\end{figure}

\begin{figure}[H]
\centering
\begin{subfigure}{0.24\textwidth}
\includegraphics[width=\textwidth]{object_c_teapot_front.png}
\caption{Front}
\end{subfigure}
\begin{subfigure}{0.24\textwidth}
\includegraphics[width=\textwidth]{object_c_teapot_side.png}
\caption{Side}
\end{subfigure}
\begin{subfigure}{0.24\textwidth}
\includegraphics[width=\textwidth]{object_c_teapot_back.png}
\caption{Back}
\end{subfigure}
\begin{subfigure}{0.24\textwidth}
\includegraphics[width=\textwidth]{object_c_teapot_side2.png}
\caption{Side 2}
\end{subfigure}
\caption{Object C: Utah Teapot (classic 3D test model). Original single-view input shown inset.}
\label{fig:object_c_views}
\end{figure}

\begin{figure}[H]
\centering
\includegraphics[width=0.85\textwidth]{render_comparison.png}
\caption{Render comparison: ground truth vs. 2DGS rendered views for Object A (Lego).}
\label{fig:render_comparison}
\end{figure}

\begin{figure}[H]
\centering
\includegraphics[width=0.75\textwidth]{fusion_keyframes.png}
\caption{Key frames from the fusion flythrough video. Objects A (left), B (center),
and C (right) are placed on the bonsai table and rendered with layered compositing.}
\label{fig:fusion_frames}
\end{figure}

The layered compositing approach successfully integrates all four elements
into a coherent scene. Minor lighting inconsistencies are visible at object
boundaries due to the different illumination conditions under which each
asset was generated. These are discussed further in Section~\ref{sec:discussion}.

% ===========================================================================
\section{Discussion and Analysis}
\label{sec:discussion}
% ===========================================================================

\subsection{Comparison of the Three Paradigms}

\begin{table}[H]
\centering
\caption{Qualitative comparison across key dimensions}
\label{tab:qualitative}
\begin{tabular}{@{}lccc@{}}
\toprule
\textbf{Dimension} & \textbf{Multi-View (2DGS)} & \textbf{Text-to-3D (threestudio)} & \textbf{Teapot Model} \\
\midrule
Geometric Accuracy      & Excellent (photometric loss) & Fair (SDS ambiguity) & Good (dual-prior constraints) \\
Texture Fidelity        & Excellent (direct photo) & Good (diffusion prior) & Good (reference anchor) \\
Computational Cost      & Low ($\sim$12 min) & High ($\sim$10 min) & Medium ($\sim$45 min) \\
Input Convenience       & Complex (100 photos + COLMAP) & Easy (1 sentence) & Easy (1 photo) \\
Output Control          & Full (explicit geometry) & Limited (prompt only) & Moderate (reference view) \\
Generalization          & Scene-specific only & Open-vocabulary & Object-agnostic \\
\bottomrule
\end{tabular}
\end{table}

\paragraph{Geometric Accuracy.}
Multi-view reconstruction achieves the highest geometric fidelity because it
directly optimizes against real photographs with explicit multi-view consistency.
Text-to-3D shows the lowest geometric accuracy due to the inherent ambiguity of
the SDS objective---the mode-seeking behavior of score distillation leads to
plausible-but-imprecise geometry. Single-image-to-3D occupies a middle ground:
the reference view provides a strong geometric anchor, while diffusion priors
complete unobserved regions with learned plausibility.

\paragraph{Texture Fidelity.}
Object A perfectly reproduces input textures---every color and pattern is
directly supervised by the photographic evidence. Object B generates textures
consistent with the text prompt but with reduced high-frequency detail (a
well-documented limitation of the SDS loss with classifier-free guidance).
Object C faithfully reproduces textures at the reference viewpoint but
degrades progressively at more extreme viewing angles.

\paragraph{Computational Cost Analysis.}
Measured training times on a single RTX 2080 Ti (11GB):
\begin{itemize}[leftmargin=*]
    \item \textbf{Object A (2DGS)}: 30,000 iterations in $\sim$12 minutes.
    The CUDA-accelerated tile-based rasterizer achieves $\sim$42 it/s at
    800$\times$800 resolution with 123K Gaussians.
    \item \textbf{Object B (threestudio/SDS)}: 500 training steps in $\sim$10 minutes
    using Stable Diffusion 1.5 with model CPU offloading to fit within 11GB VRAM.
    Each step requires forward/backward through SD UNet (860M params).
    \item \textbf{Object C (Teapot)}: Coarse (5000 iters) + Fine (3000 iters)
    stages in $\sim$45 minutes. The two-stage pipeline with Zero123 (15.5GB) +
    SD 1.5 + NeRF reaches $\sim$1.0 it/s, constrained by 11GB VRAM.
\end{itemize}

\paragraph{Practical Trade-offs.}
Multi-view reconstruction requires the most capture effort but delivers the
highest quality---suitable for applications demanding photorealistic accuracy
(e.g., digital heritage, product visualization). Text-to-3D requires only a
prompt, making it ideal for rapid prototyping and creative exploration, but
offers the least control over output. Single-image-to-3D represents the most
balanced trade-off, providing good quality from minimal input---practical for
many real-world scenarios.

\subsection{Unified Representation for Heterogeneous 3D Assets}
\label{sec:unified}

A core technical challenge is that the background (bonsai) and Object A (Lego)
are represented as 2D Gaussian primitives, while Objects B and C are textured
triangle meshes from threestudio (NeRF$\rightarrow$marching cubes) and Magic123
(DMTet). We adopt a \textbf{layered rendering with vertex-color mesh
rasterization} approach:

\begin{enumerate}[leftmargin=*]
    \item \textbf{2DGS layers (Background + Object A)}: Rendered natively using
    the CUDA-accelerated 2DGS differentiable rasterizer. Object A's Gaussian
    coordinates are transformed (scaled $\times$0.85, translated) into the bonsai
    COLMAP coordinate frame.
    \item \textbf{Mesh layers (Objects B + C)}: Rendered using Open3D's
    Filament-based OffscreenRenderer with the \texttt{defaultUnlit} shader.
    Vertex colors are sampled from UV texture maps (texture\_kd.jpg for B,
    albedo.png for C) at each vertex's UV coordinate using bilinear
    interpolation: $C_i = T(u_i, v_i)$ where $T$ is the texture image.
    This approach was chosen over Gaussian surface sampling because:
    \begin{itemize}
        \item Preserves sharp mesh geometry without Gaussian approximation
        \item Directly uses the mesh's UV texture mapping for accurate colors
        \item Avoids the memory-intensive Gaussian rasterization of dense
        surface samples
    \end{itemize}
    \item \textbf{Alpha compositing}: Each layer produces an RGBA image:
    \begin{itemize}
        \item Background: rendered with black background ($\alpha=1$)
        \item Object A: rendered with white background; $\alpha = \max(1-R, 1-G, 1-B) \times 3.0$ extracts the object silhouette
        \item Objects B/C: rendered with transparent background; $\alpha = \mathbb{I}(\max(R,G,B) > 0.01)$
    \end{itemize}
    Layers are composited front-to-back: $\text{result} = C_A\alpha_A + C_{\text{bg}}(1-\alpha_A)$, etc.
    \item \textbf{Camera system unification}: All layers share the same virtual
    camera defined by world-space eye position, look-at point, and vertical field
    of view (50$^\circ$). The 2DGS camera and Open3D camera are independently
    configured to match these parameters, ensuring pixel-perfect alignment.
\end{enumerate}

This approach achieves \textbf{code-level fusion} without requiring mesh-to-Gaussian
conversion or external tools like Blender. All rendering is performed
programmatically, enabling a fully automated pipeline from training to video.

\subsection{Fusion Challenges and Limitations}

\begin{enumerate}[leftmargin=*]
    \item \textbf{Lighting inconsistency}: Each asset was generated under different
    virtual illumination conditions (white studio lighting for Object A,
    diffusion-prior ambient for B, image-based for C, outdoor daylight for
    background). The composited scene shows visible lighting seams. Future work
    should incorporate relighting techniques or physically-based environment
    mapping.

    \item \textbf{Scale calibration}: Automatic scale estimation from COLMAP
    reconstructions is unreliable for monocular SfM. We manually adjusted object
    scales based on scene context, but automated methods (e.g., using known
    reference objects or learned scale priors) would improve robustness.

    \item \textbf{Representation heterogeneity}: Converting 2DGS (explicit
    primitives) to mesh format for nvdiffrast rendering introduces approximation
    error. The Gaussian-to-mesh conversion via Poisson surface reconstruction
    loses the view-dependent appearance modeling capability of the original
    representation.

    \item \textbf{Occlusion handling}: Our depth-based compositing assumes
    per-object mean depth ordering, which fails when objects overlap in depth
    from certain viewpoints. Per-pixel depth compositing would be more accurate
    but requires depth map rendering for all components.

    \item \textbf{Missing shadows and inter-reflections}: Objects are composited
    without casting shadows on the background or on each other, which reduces
    realism. Physically-based rendering with environment maps would address this.
\end{enumerate}

% ===========================================================================
\section{Conclusion}
% ===========================================================================

We have implemented and systematically compared three contemporary paradigms for
3D asset creation within a unified experimental framework. Our findings confirm
that multi-view photogrammetry (2DGS) achieves the highest geometric and
photometric quality at the cost of extensive data capture requirements; text-to-3D
generation (threestudio/SDS) offers the greatest creative flexibility with
minimal input effort; and single-image reconstruction (Magic123) provides the
best quality-to-convenience ratio for practical applications.

The fusion pipeline demonstrates that heterogeneous 3D representations can be
successfully composited through layer-based image-space rendering, enabling the
creation of rich scenes that combine real scanned environments with AI-generated
virtual assets. While limitations remain in lighting consistency and shadow
modeling, the complete pipeline provides a practical foundation for full-stack
3D vision applications.

Future work directions include: (1) incorporating differentiable path tracing for
physically-accurate compositing; (2) automated scale and pose alignment using
learned correspondence; (3) extending text-to-3D with ControlNet-based geometric
guidance; and (4) real-time fusion rendering for interactive applications.

% ===========================================================================
\section*{Project Links}
% ===========================================================================

\noindent
\begin{tabular}{@{}ll@{}}
\textbf{GitHub Repository:} & \url{https://github.com/YOUR_USERNAME/3d-vision-pipeline} \\
& Extraction code: \texttt{xxxx} \\
\end{tabular}

\vspace{8pt}
\noindent\textit{Note: Replace placeholder URLs above with actual public links before submission.}

% ===========================================================================
\bibliographystyle{unsrt}
\begin{thebibliography}{99}

\bibitem{huang20242dgs}
B.~Huang, Z.~Yu, A.~Chen, A.~Geiger, and S.~Gao,
``2D Gaussian Splatting for Geometrically Accurate Radiance Fields,''
in \emph{ACM SIGGRAPH}, 2024.

\bibitem{kerbl20233dgs}
B.~Kerbl, G.~Kopanas, T.~Leimkühler, and G.~Drettakis,
``3D Gaussian Splatting for Real-Time Radiance Field Rendering,''
in \emph{ACM SIGGRAPH}, 2023.

\bibitem{schoenberger2016sfm}
J.~L.~Schönberger and J.-M.~Frahm,
``Structure-from-Motion Revisited,''
in \emph{CVPR}, 2016.

\bibitem{poole2022dreamfusion}
B.~Poole, A.~Jain, J.~T.~Barron, and B.~Mildenhall,
``DreamFusion: Text-to-3D using 2D Diffusion,''
in \emph{ICLR}, 2023.

\bibitem{guo2023threestudio}
Y.-C.~Guo et al.,
``threestudio: A unified framework for 3D content generation,''
\url{https://github.com/threestudio-project/threestudio}, 2023.

\bibitem{qian2023magic123}
G.~Qian et al.,
``Magic123: One Image to High-Quality 3D Object Generation Using Both 2D and 3D
Diffusion Priors,''
in \emph{CVPR}, 2024.

\bibitem{mildenhall2020nerf}
B.~Mildenhall, P.~P.~Srinivasan, M.~Tancik, J.~T.~Barron, R.~Ramamoorthi, and R.~Ng,
``NeRF: Representing Scenes as Neural Radiance Fields for View Synthesis,''
in \emph{ECCV}, 2020.

\bibitem{barron2022mipnerf360}
J.~T.~Barron, B.~Mildenhall, D.~Verbin, P.~P.~Srinivasan, and P.~Hedman,
``Mip-NeRF 360: Unbounded Anti-Aliased Neural Radiance Fields,''
in \emph{CVPR}, 2022.

\bibitem{laine2020nvdiffrast}
S.~Laine, J.~Hellsten, T.~Karras, Y.~Seol, J.~Lehtinen, and T.~Aila,
``Modular Primitives for High-Performance Differentiable Rendering,''
in \emph{ACM SIGGRAPH Asia}, 2020.

\bibitem{rombach2022high}
R.~Rombach, A.~Blattmann, D.~Lorenz, P.~Esser, and B.~Ommer,
``High-Resolution Image Synthesis with Latent Diffusion Models,''
in \emph{CVPR}, 2022.

\end{thebibliography}

\end{document}
